\documentclass[review,11pt,3p,authoryear,nopreprintline]{elsarticle}
\usepackage[T1]{fontenc}
\usepackage{amsmath,amssymb,bm}
\pagestyle{empty}
\begin{document}
\noindent\textbf{Figures 1 and 2: actual manuscript inclusion sizes}

\noindent Body-text reference (10.95 pt): empirical probability $q$, candidate
probability $p$, response threshold $c$, and probability $p_f(c)$.
The main figure lettering matches this size; Figure 1's inset details use 9.5 pt.

\noindent Figure 1: $0.65\,\mathtt{columnwidth}$.
\begin{center}
  \includegraphics[width=0.65\columnwidth]{figures/kl_dw_comparison/fig_kl_dw_ambiguity_comparison.pdf}
\end{center}
\noindent Figure 2: $0.50\,\mathtt{columnwidth}$.
\begin{center}
  \includegraphics[width=0.50\columnwidth]{figures/response_threshold_band_optimized/fig_response_threshold_band.pdf}
\end{center}
\clearpage
\noindent\textbf{Figures 3 and 4: actual manuscript inclusion sizes}

\noindent Body-text reference (10.95 pt): residual cells $\Xi_j$, input cells
$\mathcal X_i$, probabilities $w_i$, and joint cell masses $\pi_{ij}$.

\noindent Figure 3: $0.60\,\mathtt{columnwidth}$.
\begin{center}
  \includegraphics[width=0.60\columnwidth]{figures/residual_input_cells_optimized/fig_residual_input_cells.pdf}
\end{center}
\noindent Figure 4: $0.60\,\mathtt{columnwidth}$.
\begin{center}
  \includegraphics[width=0.60\columnwidth]{figures/joint_partition_optimized/fig_joint_partition.pdf}
\end{center}
\end{document}
